% !TEX program = xelatex
% ------------------------------------------------------------------------
% main.tex —— TNDQ 论文主文件
% 编译：xelatex main.tex（建议连续编译两遍以生成正确的交叉引用）
% ------------------------------------------------------------------------
\documentclass[11pt,a4paper]{article}

\usepackage{tndq-style}

\title{Third-Order Dual Quaternion (TODQ) Kinematics with Geometrically
       Consistent Error Systems:\\
       A Computed Torque Approach to $H_\infty$/ISS Control}
\author{（作者姓名）\\（单位，城市，国家）}
\date{}

\begin{document}

\maketitle
\thispagestyle{fancy}

% ================================ 摘要 ================================
\begin{abstract}
\noindent
对偶四元数（DQ）为机械臂位姿提供了全局无奇异参数化，但现有框架仅携带零阶信息，向动力学接口延伸时存在位姿/速度误差分离、加速度扰动无入口、偏差型不确定性不满足 $L_2$ 假设三个结构性缺口。本文引入\textbf{三阶对偶四元数}（TODQ）代数 $\mathcal A_2$ \cite{Con25}，使正运动学一次连乘同时输出位姿、twist 与任务空间加速度。在此基础上，利用 HDQ 截断与乘法的相容性，将误差体系定义在两项 HDQ 上：一次乘法同时生成右不变位姿误差与几何一致 twist 误差（定理 1），并导出严格级联误差运动学（定理 2）。针对动力学接口设计几何一致计算力矩律，证明闭环满足精确耗散等式 $\dot V=-e_\xi^\top K_de_\xi+e_\xi^\top d$（无任何放缩），据此给出无扰时的水平集不变性与渐近收敛、$L_2$ 扰动下 H∞ 增益的 Schur 补充分必要条件，以及 $L_\infty$ 扰动下 twist 误差的均方极限界（定理 3）。在 7 自由度机械臂力矩模式仿真中验证静态刚度标度律与定理 3 的性能保证。
\end{abstract}

\medskip
\noindent\textbf{关键词}：对偶四元数；超对偶四元数；三阶对偶四元数；多项式代数；几何一致误差；H∞ 控制；输入-状态稳定

\bigskip

% ================================ 正文 ================================
% ---------------------------------------------------------------------
% 第 1 章 引言
% ---------------------------------------------------------------------
\section{引言}

\subsection{背景与动机}

刚体位姿（姿态 + 位置）的参数化是机器人控制的起点。单位对偶四元数将两者装入一个 8 维代数对象，运算全局无奇异，且乘法直接实现位姿复合；\cite{P2} 在此参数化上给出了带 $L_2$ 扰动衰减保证的运动学跟踪控制器，是 DQ 控制的代表性结果。Cohen \& Shoham \cite{P1} 进一步引入超对偶四元数（HDQ）：给 DQ 增加一个幂零单位 $\varepsilon^*$（$\varepsilon^{*2}=0$），使 HDQ 能同时表示位姿与其一阶时间导数，链式乘法自动微分，从而正运动学一次连乘同时输出位姿与 twist。

但 HDQ 到动力学层面仍差一阶：计算力矩控制需要任务空间加速度 $\dot{\boldsymbol\xi}$ 与雅可比导数项 $\dot J\dot{\boldsymbol q}$，而 \cite{P1} 的 HDQ 结构中两个幂零单位（$\varepsilon$ 承载平移、$\varepsilon^*$ 承载一阶导）均已占用，$\varepsilon\varepsilon^*$ 项只表示“平移分量的一阶导数”，不含二阶时间导数。更重要的是\textbf{误差体系}：现有 DQ 框架的误差对象只有位姿一层，扰动模型假设速度级加性扰动且属于 $L_2$；进入动力学接口后出现三个结构性缺口，即位姿与速度误差分离（\cite{Ch20} 类律的位姿反馈取螺旋对数，其导数映射在大误差处奇异）、加速度层扰动无入口、偏差型不确定性不满足 $L_2$ 假设，既有证书随之失效。

\subsection{本文贡献}

\begin{enumerate}[leftmargin=2em]
  \item \textbf{TODQ 运动学}：基于三项多项式代数构造三阶对偶四元数串联运动学，突破现有 HDQ 正运动学止于一阶导数的限制，一次 $O(n)$ 连乘同时输出位姿、twist 与任务空间加速度，为加速度级控制提供解析前馈。
  \item \textbf{一体化误差体系}：在两项 HDQ 元素上定义误差对象，一次乘法同时生成右不变位姿误差与几何一致 twist 误差，并证明闭环误差动力学为严格级联标准形。
  \item \textbf{控制律与精确证书}：提出 $A^\top$ 整形计算力矩律，闭环存储函数满足精确耗散等式，统一给出无扰收敛、$L_2$ 衰减充要判据与偏差型扰动的均方界；$L_2$ 增益经最优化恰为最小阻尼的倒数，与线性化真实 H∞ 范数相等，衰减指标可直接换算为阻尼下界。
  \item \textbf{实机验证}：在六自由度机械臂上完成四律同预算对比与 $\gamma$ 增益阶梯共 33 次运行，各证书全程成立，并区分任务律控制的瞬态精度与拴定弹簧、摩擦主导的准静态残差。
\end{enumerate}

TODQ 是 \cite{P1} HDQ（幂零单位承载导数）向二阶的最小扩展 \cite{Con25}；0 阶误差退化为 \cite{P2}（§4.4）。

% ---------------------------------------------------------------------
% 第 2 章 预备知识
% ---------------------------------------------------------------------
\section{预备知识}

\textbf{记号约定}（全文统一）：同一条位姿曲线用同一核心字母（如 $x$），字母上方的装饰指明它所处的代数层：

\begin{center}
\begin{tabularx}{\linewidth}{@{}llX@{}}
\toprule
记号 & 对象 & 说明 \\
\midrule
$\hat a$ & 四元数 & $\hat a\in\mathrm{Spin}(3)$，§2.1 \\
$\hat{\underline a}$ & 对偶四元数 & $\hat{\underline a}\hat{\underline a}^*=1$，式 \eqref{eq:2.1} \\
$\breve a$ & HDQ（超对偶四元数） & 两个 DQ 项；曲线的 HDQ 表示 $\breve x=\hat{\underline x}+\varepsilon^*\dot{\hat{\underline x}}$，§2.3 \\
$\bar a$ & TODQ（三阶对偶四元数） & 三个 DQ 项；曲线的 TODQ 表示 $\bar x=\hat{\underline x}+\sigma\dot{\hat{\underline x}}+\tfrac12\sigma^2\ddot{\hat{\underline x}}$，§3.2 \\
$\tilde{(\cdot)}$ & 误差量 & 只佩戴波浪号，不再叠加类型装饰；其类型由定义式指明（如 $\tilde x=\hat{\underline x}\hat{\underline x}_d^{\,*}$ 是单位 DQ，$\tilde r$ 是单位四元数） \\
粗体 & 纯 DQ（twist 等）、向量与矩阵 & $\boldsymbol\xi,\boldsymbol q,J,K_d$ 等 \\
\bottomrule
\end{tabularx}
\end{center}

（表 0：记号约定。标称模型矩阵 $\hat M,\hat C,\hat g$ 上的 hat 沿用控制文献“标称估计”的习惯用法，与四元数装饰无关；四元数虚单位 $\hat\imath,\hat\jmath,\hat k$ 为固定符号；对偶四元数代数统一记作 DQ（§2.1），其元素是一般（未必单位）对偶四元数；带下标的项记号 $\hat{\underline a}_k,\hat{\underline b}_k$（§2.3、定义 1）表示一般对偶四元数项，未必单位——如曲线表示的导数项 $\dot{\hat{\underline x}}$。）

\subsection{四元数与对偶四元数}

单位四元数 $\hat r=\cos\frac\phi2+n\sin\frac\phi2$ 表示绕单位轴 $n$ 转角 $\phi$ 的旋转。

对偶四元数（DQ）代数为 $\mathbb H\oplus\varepsilon\mathbb H$，$\varepsilon^2=0$（$\varepsilon$ 与四元数单位交换）。\textbf{单位 DQ}
\begin{equation}
\hat{\underline x}=\hat r+\varepsilon\tfrac12\,p\,\hat r,\qquad \hat r\in\mathrm{Spin}(3),\ p\in\mathbb H_p\ (\text{纯四元数}\cong\mathbb R^3)
\tag{2.1}\label{eq:2.1}
\end{equation}
表示位姿（旋转 $\hat r$ + 平移 $p$），满足 $\hat{\underline x}\hat{\underline x}^*=1$；单位 DQ 全体构成群 $\mathrm{Spin}(3)\ltimes\mathbb R^3$，双覆盖 $SE(3)$，乘法即位姿复合。DQ 共轭 $\hat{\underline x}^*=\hat r^*+\varepsilon\tfrac12\hat r^*p^*$ 逐分量遵循四元数共轭。\textbf{纯 DQ}（标量部与对偶标量部为零者）构成 6 维实空间，记 $\mathrm{vec}_6$ 为取其两个向量部的坐标同构，$\overline{\mathrm{vec}}_6$ 为逆。

\subsection{twist 与运动学}

沿光滑单位 DQ 曲线 $\hat{\underline x}(t)$，\textbf{空间 twist} 定义为
\begin{equation}
\boldsymbol\xi\triangleq 2\dot{\hat{\underline x}}\hat{\underline x}^{*}=\omega+\varepsilon v,\qquad v=\dot p+p\times\omega,
\tag{2.2}\label{eq:2.2}
\end{equation}
其中 $\omega$ 为角速度。$\boldsymbol\xi$ 是纯 DQ，等价写法即左乘运动学 $\dot{\hat{\underline x}}=\tfrac12\boldsymbol\xi\hat{\underline x}$（\cite{P2} 式(1) 约定）。对 $n$ 关节串联机械臂，正运动学 $\hat{\underline x}(\boldsymbol q)=\prod_{i=1}^n\hat{\underline x}_i(q_i)$（各关节单位 DQ 之积），微分给出 $\mathrm{vec}_6\,\boldsymbol\xi=J(\boldsymbol q)\dot{\boldsymbol q}$，$J\in\mathbb R^{6\times n}$ 称几何雅可比。

\textbf{伴随作用与李括号}：$\mathrm{Ad}_{\hat{\underline x}}\boldsymbol a\triangleq\hat{\underline x}\boldsymbol a\hat{\underline x}^*$（twist 的参考系搬运）、$\mathrm{ad}_{\boldsymbol a}\boldsymbol b\triangleq\tfrac12(\boldsymbol{ab}-\boldsymbol{ba})$（李括号的 DQ 形式），二者均保持纯 DQ 结构，详见 \cite{Lynch17}。

\subsection{超对偶四元数（HDQ）}

HDQ 元素形如 $\breve a=\hat{\underline a}_0+\varepsilon^*\hat{\underline a}_1$，两个项 $\hat{\underline a}_0,\hat{\underline a}_1$ 均为 DQ（§2.1）；$\varepsilon^*$ 与全部四元数单位交换且 $\varepsilon^{*2}=0$，乘法由分配律给出：
\begin{equation}
(\hat{\underline a}_0+\varepsilon^*\hat{\underline a}_1)(\hat{\underline b}_0+\varepsilon^*\hat{\underline b}_1)=\hat{\underline a}_0\hat{\underline b}_0+\varepsilon^*(\hat{\underline a}_0\hat{\underline b}_1+\hat{\underline a}_1\hat{\underline b}_0).
\tag{2.3}\label{eq:2.3}
\end{equation}

\textbf{HDQ 共轭}（定理 1 中 $(\breve x_d)^*$ 的含义）定义为\textbf{逐项的 DQ 共轭}：
\begin{equation}
\breve a^{\,*}\triangleq\hat{\underline a}_0^{\,*}+\varepsilon^*\hat{\underline a}_1^{\,*}.
\tag{2.4}\label{eq:2.4}
\end{equation}
\eqref{eq:2.4} 使共轭与求导可交换（$ (\breve x)^* $ 即曲线 $\hat{\underline x}^*(t)$ 的 HDQ 表示），且仍为反自同构：$ (\breve a\breve b)^*=\breve b^{\,*}\breve a^{\,*} $。

% ---------------------------------------------------------------------
% 第 3 章 TODQ：三阶对偶四元数与运动学重构
% ---------------------------------------------------------------------
\section{TODQ：三阶对偶四元数与运动学重构}

\subsection{代数定义}

\begin{definition}[TODQ \cite{Con25}]\label{def:todq}
三阶对偶四元数（TODQ）的元素形如
\begin{equation}
\bar a=\hat{\underline a}_0+\sigma\hat{\underline a}_1+\tfrac12\sigma^2\hat{\underline a}_2,\qquad \hat{\underline a}_0,\hat{\underline a}_1,\hat{\underline a}_2\ \text{均为 DQ}.
\tag{3.1}\label{eq:3.1}
\end{equation}
全部 TODQ 元素构成的代数记 $\mathcal A_2$。$\sigma$ 与全部四元数单位交换，$\sigma^3=0$。乘法由分配律与 $\sigma^3=0$ 唯一确定：
\begin{equation}
\bar a\,\bar b=\hat{\underline a}_0\hat{\underline b}_0+\sigma(\hat{\underline a}_0\hat{\underline b}_1+\hat{\underline a}_1\hat{\underline b}_0)+\tfrac12\sigma^2\bigl(\hat{\underline a}_0\hat{\underline b}_2+2\hat{\underline a}_1\hat{\underline b}_1+\hat{\underline a}_2\hat{\underline b}_0\bigr).
\tag{3.2}\label{eq:3.2}
\end{equation}
\textbf{TODQ 共轭}定义为逐项 DQ 共轭（与 \eqref{eq:2.4} 同一约定）：$\bar a^{\,*}\triangleq\hat{\underline a}_0^{\,*}+\sigma\hat{\underline a}_1^{\,*}+\tfrac12\sigma^2\hat{\underline a}_2^{\,*}$，它是 $\mathcal A_2$ 上的反自同构 $(\bar a\bar b)^*=\bar b^{\,*}\bar a^{\,*}$。
\end{definition}

\subsection{TODQ 串联运动学表示}

给定光滑单位 DQ 曲线 $\hat{\underline x}(t)$，其\textbf{TODQ 表示}定义为把 $\hat{\underline x}$ 与它的两阶导数按 $\sigma$ 的幂次装入三个项：
\begin{equation}
\bar x\triangleq\hat{\underline x}+\sigma\dot{\hat{\underline x}}+\tfrac12\sigma^2\ddot{\hat{\underline x}}\ \in\mathcal A_2 .
\tag{3.3a}\label{eq:3.3a}
\end{equation}

对串联臂 $\hat{\underline x}(\boldsymbol q)=\prod_i\hat{\underline x}_i(q_i(t))$，先写出每个关节因子的 TODQ 表示 $\bar x_i$（由 $\partial\hat{\underline x}_i/\partial q_i=\tfrac12\boldsymbol s_i\hat{\underline x}_i$ 与链式法则给出闭式，只依赖 $q_i,\dot q_i,\ddot q_i$），再按 \eqref{eq:3.2} 逐个连乘：
\begin{equation}
\bar x=\bar x_1\,\bar x_2\cdots\bar x_n=\prod_{i=1}^{n}\bar x_i
\tag{3.4}\label{eq:3.4}
\end{equation}
一次 $O(n)$ 链连乘同时输出 $\hat{\underline x},\dot{\hat{\underline x}},\ddot{\hat{\underline x}}$（$\bar x$ 的三个项）。导出量：
\begin{equation}
\boldsymbol\xi=2\dot{\hat{\underline x}}\hat{\underline x}^*,\qquad
\dot{\boldsymbol\xi}=2\ddot{\hat{\underline x}}\hat{\underline x}^*-\tfrac12\boldsymbol\xi^2\ \text{（取纯部）},\qquad
\mathrm{vec}_6\dot{\boldsymbol\xi}=\dot J\dot{\boldsymbol q}+J\ddot{\boldsymbol q}.
\tag{3.5}\label{eq:3.5}
\end{equation}

% ---------------------------------------------------------------------
% 第 4 章 几何一致误差体系
% ---------------------------------------------------------------------
\section{几何一致误差体系}

\subsection{误差对象的阶数}

计算力矩律的反馈项为 $-K_de_\xi-A^\top K_pe_z$，只需位姿与速度误差；期望加速度用于前馈，动力学失配归入 $\dot e_\xi$ 中的扰动。因此正运动学与期望轨迹用 TODQ 建模，误差体系定义在 HDQ 上，无需测量加速度误差。

\subsection{定理 1：误差的 HDQ 表示}

实测链取 HDQ 表示 $\breve x=\hat{\underline x}+\varepsilon^*\dot{\hat{\underline x}}$；期望链取 $\breve x_d=\hat{\underline x}_d+\varepsilon^*\dot{\hat{\underline x}}_d$。

\begin{theorem}[误差的 HDQ 提升]\label{thm:1}
本文采用右误差 $\tilde x=x x_d^*$，并使用空间表示的 twist。定义 HDQ 误差元素
\begin{equation}
\breve{\tilde x}\triangleq \breve x\cdot\bigl(\breve x_d\bigr)^{*} .
\tag{4.1}\label{eq:4.1}
\end{equation}
则
\begin{equation}
\breve{\tilde x}=\tilde x+\varepsilon^*\dot{\tilde x},\qquad
\tilde x=\hat{\underline x}\hat{\underline x}_d^{\,*},\quad
\dot{\tilde x}=\dot{\hat{\underline x}}\hat{\underline x}_d^{\,*}+\hat{\underline x}\dot{\hat{\underline x}}_d^{\,*},
\tag{4.2}\label{eq:4.2}
\end{equation}
即一次 HDQ 乘法（3 次 DQ 乘）同时给出位姿误差 $\tilde x$ 与其导数；0 阶项正是 \cite{P2} 的 $\tilde x$。进一步定义\textbf{几何一致 twist 误差}
\begin{equation}
\tilde{\boldsymbol\xi}\triangleq 2\dot{\tilde x}\tilde x^{*},\qquad
e_\xi\triangleq\mathrm{vec}_6\tilde{\boldsymbol\xi},
\tag{4.3}\label{eq:4.3}
\end{equation}
则：(i) $\tilde{\boldsymbol\xi}$ 是纯 DQ，且对 unwinding 翻转 $\tilde x\to-\tilde x$ 不变；(ii) 误差满足与原系统同型的左乘运动学 $\dot{\tilde x}=\tfrac12\tilde{\boldsymbol\xi}\tilde x$；(iii) 由误差定义和运动学关系有
\begin{equation}
\tilde{\boldsymbol\xi}=\boldsymbol\xi-\mathrm{Ad}_{\tilde x}\boldsymbol\xi_d ,
\tag{4.4}\label{eq:4.4}
\end{equation}
即“当前 twist 减去\textbf{搬运到当前位姿处}的期望 twist”。
\end{theorem}

\begin{proof}
式 \eqref{eq:4.2} 由 HDQ 乘法与 Leibniz 法则直接得到。(i) 对 $\tilde x\tilde x^*=1$ 求导可得 $\tilde{\boldsymbol\xi}+\tilde{\boldsymbol\xi}^*=0$，故其为纯 DQ；同时翻转 $\tilde x,\dot{\tilde x}$ 不改变 $2\dot{\tilde x}\tilde x^*$。(ii) 式 \eqref{eq:4.3} 右乘 $\tilde x$ 即得。(iii) 由 $\dot{\hat{\underline x}}_d^{\,*}=-\tfrac12\hat{\underline x}_d^{\,*}\boldsymbol\xi_d$ 得 $\dot{\tilde x}=\tfrac12\boldsymbol\xi\tilde x-\tfrac12\tilde x\boldsymbol\xi_d$，右乘 $2\tilde x^*$ 即 \eqref{eq:4.4}。
\end{proof}

\subsection{定理 2：输出误差运动学}

沿用 \cite{P2} 输出 $e_z=[\mathcal O;\mathcal T]\in\mathbb R^6$，$\mathcal O=-\mathrm{Im}\,\tilde r$，$\mathcal T=\tilde p$（$\tilde r,\tilde p$ 为 $\tilde x$ 的旋转/归一化平移分量，$\tilde r=\tilde\eta+\tilde\mu$）。

\begin{theorem}[输出误差运动学闭式]\label{thm:2}
记 $\tilde{\boldsymbol\xi}=\tilde\omega+\varepsilon\tilde v$，则
\begin{equation}
\dot e_z=A(\tilde x)\,e_\xi,
\qquad
A(\tilde x)=
\begin{bmatrix}
-\tfrac12\bigl(\tilde\eta I_3+[\mathcal O]_\times\bigr) & 0_3\\[2pt]
-[\mathcal T]_\times & I_3
\end{bmatrix},
\tag{4.5}\label{eq:4.5}
\end{equation}
且 $\tilde x\to1$ 时 $A\to A_0=\mathrm{diag}(-\tfrac12I_3,I_3)$，$\sigma_{\min}(A_0)=\tfrac12$。
\end{theorem}

\begin{proof}
旋转分量：由 $\dot{\tilde r}=\tfrac12\tilde\omega\tilde r$（定理 1(ii) 的四元数部），$\dot{\mathcal O}=-\mathrm{Im}(\tfrac12\tilde\omega\tilde r)=-\tfrac12(\tilde\eta I_3+[\mathcal O]_\times)\tilde\omega$（用 $\tilde\omega\times\tilde\mu=[\mathcal O]_\times\tilde\omega$）。平移分量：由 twist 约定 $\tilde v=\dot{\tilde p}+\tilde p\times\tilde\omega$ 得 $\dot{\mathcal T}=\tilde v-[\mathcal T]_\times\tilde\omega$。
\end{proof}

定理 1、2 给出二阶误差系统：$\dot e_z=Ae_\xi$，$\dot e_\xi$ 由控制律与扰动决定（§5），无需另设加速度误差状态。

% ---------------------------------------------------------------------
% 第 5 章 几何一致控制律与认证 H∞ 性能
% ---------------------------------------------------------------------
\section{几何一致控制律与认证 H∞ 性能}

\subsection{扰动项：显式解算、乘性项分离与假设集}

真实关节空间动力学（含摩擦 $\boldsymbol f$、执行器实现误差 $\delta\boldsymbol\tau$ 与环境接触 $\boldsymbol\tau_{\mathrm{ext}}$）记为 $M(\boldsymbol q)\ddot{\boldsymbol q}+C(\boldsymbol q,\dot{\boldsymbol q})\dot{\boldsymbol q}+\boldsymbol g(\boldsymbol q)+\boldsymbol f=\boldsymbol\tau+\delta\boldsymbol\tau+\boldsymbol\tau_{\mathrm{ext}}$；内环以标称模型执行计算力矩 $\boldsymbol\tau=\hat M\ddot{\boldsymbol q}_{\mathrm{ref}}+\hat C\dot{\boldsymbol q}+\hat g+\hat{\boldsymbol f}$。失配量取“标称减真实”：$\Delta M\triangleq\hat M-M$，$\Delta C,\Delta\boldsymbol g,\Delta\boldsymbol f$ 同理。代入并左乘 $M^{-1}$ 得 $\ddot{\boldsymbol q}=\ddot{\boldsymbol q}_{\mathrm{ref}}+\boldsymbol w_{\mathrm{dyn}}$，其中
\begin{equation}
\boldsymbol w_{\mathrm{dyn}}=M^{-1}\bigl(\Delta M\,\ddot{\boldsymbol q}_{\mathrm{ref}}+\Delta C\dot{\boldsymbol q}+\Delta\boldsymbol g+\Delta\boldsymbol f+\delta\boldsymbol\tau+\boldsymbol\tau_{\mathrm{ext}}\bigr).
\tag{5.1}\label{eq:5.1}
\end{equation}

\textbf{反馈乘性项与余项分离}。记
\[
u_{\rm ff}\triangleq\mathrm{vec}_6\!\left(\mathrm{Ad}_{\tilde x}\dot{\boldsymbol\xi}_d+\mathrm{ad}_{\tilde{\boldsymbol\xi}}\mathrm{Ad}_{\tilde x}\boldsymbol\xi_d\right),\qquad
u_{\rm fb}\triangleq-K_de_\xi-A^{\top}(\tilde x)K_pe_z .
\]
将 $\ddot{\boldsymbol q}_{\mathrm{ref}}=J^{+}(u_{\rm ff}+u_{\rm fb}-\dot J\dot{\boldsymbol q})$ 代入 \eqref{eq:5.1}，利用 $JJ^{+}=I_6$ 得扰动分解
\begin{equation}
d=\Theta(\boldsymbol q)\,u_{\mathrm{fb}}+d_{\mathrm{ex}},\qquad
\Theta\triangleq J M^{-1}\Delta M\,J^{+}\in\mathbb R^{6\times6},\;
\tag{5.1d}\label{eq:5.1d}
\end{equation}
其中 $\Theta$ 为反馈乘性分量，$d_{\mathrm{ex}}$ 为含状态相关项的余项（附录 C.1）。定理 3(c) 假设总扰动 $d\in L_2$，3(d) 则假设余项有界。

\textbf{假设集}（各结论按需引用）：
\begin{itemize}[leftmargin=2em]
  \item \textbf{(A1) 雅可比正则}：在工作集 $\mathcal Q$ 上 $\sigma_{\min}(J)\ge\underline\sigma>0$，理论取 Moore--Penrose 伪逆 $J^{+}$，故 $JJ^{+}=I_6$。阻尼伪逆产生的 $JJ^{\#}-I_6$ 残差不在此精确实现模型内。
  \item \textbf{(A2) 惯量与失配有界}：$M(\boldsymbol q)\succeq \underline m I_n$（$\underline m>0$），且 $\Delta M,\Delta C,\Delta\boldsymbol g,\Delta\boldsymbol f$ 在 $\mathcal Q\times\{\|\dot{\boldsymbol q}\|\le\bar v\}$ 上有界。
  \item \textbf{(A3) 乘性小增益条件}：
  \begin{equation}
  \alpha\triangleq\sup_{\boldsymbol q\in\mathcal Q}\bigl\|\Theta(\boldsymbol q)\bigr\|_2
  \;<\;\frac{\lambda_{\min}(K_d)}{\lambda_{\max}(K_d)}=\frac1{\mathrm{cond}_2(K_d)}\ \le 1 .
  \tag{5.1f}\label{eq:5.1f}
  \end{equation}
  各向同性阻尼时退化为 $\alpha<1$。
  \item \textbf{(A4) 轨迹与扰动正则性}：$\hat{\underline x}_d(t)$ 为 $C^2$ 且 $\|\boldsymbol\xi_d\|,\|\dot{\boldsymbol\xi}_d\|$ 有界。3(c) 要求总扰动 $d\in L_2$；3(d) 要求 $d_{\mathrm{ex}}\in L_\infty$，二者为不同情形。
\end{itemize}

\subsection{控制律}

取 $K_p$ 对称正定（典型取块对角 $K_p=\mathrm{diag}(p_O I_3,p_T I_3)$；取 $K_p=k_pI_6$ 即恢复单标量增益情形）、$K_d$ 对称正定（若需定理 3(c-2) 的旋转/平移逐分量 H∞ 界，则取块对角 $K_d=\mathrm{diag}(K_\omega,K_v)$，$K_\omega,K_v\in\mathbb R^{3\times3}$ 对称正定），定义参考加速度（几何一致计算力矩律）：
\begin{equation}
\ddot{\boldsymbol q}_{\mathrm{ref}}
=J^{+}\Bigl(\underbrace{\mathrm{vec}_6\bigl(\mathrm{Ad}_{\tilde x}\dot{\boldsymbol\xi}_d+\mathrm{ad}_{\tilde{\boldsymbol\xi}}\mathrm{Ad}_{\tilde x}\boldsymbol\xi_d\bigr)}_{\text{前馈：搬运的期望加速度 + 输运修正}}
-K_d\,e_\xi-A^{\top}(\tilde x)\,K_p\,e_z-\dot J\dot{\boldsymbol q}\Bigr).
\tag{5.2}\label{eq:5.2}
\end{equation}
其中 $A^\top$-整形位姿反馈使 Lyapunov 导数的交叉项精确相消（定理 3 证明），前馈与 $\dot J\dot{\boldsymbol q}$ 分别取自期望链的 $\sigma^2$ 项与 \eqref{eq:3.5}；全部反馈量取自定理 1 的 HDQ 误差元素，不需要任何加速度误差的测量或估计。

\subsection{主定理}

以下分析理想连续误差模型。机械臂实现另假设解在所考察时段内存在并留在 (A1)、(A2) 的正则工作集，期望轨迹满足 (A4)；无穷时域结论要求该前提全时成立，不能由任务误差有界推出。四元数采用连续提升，仅在初始时选符号。统一存储函数为
\begin{equation}
V(e_z,e_\xi)=\tfrac12\|e_\xi\|^2+\tfrac12\,e_z^{\top}K_pe_z\;\ge0 ,
\tag{5.4a}\label{eq:5.4a}
\end{equation}
$K_p\succ0$ 对称。无扰收敛限于正四元数分支的水平集；耗散不等式不要求 $A$ 可逆。

% ---- 定理 3(a)：手工编号 ----
\setcounter{theorem}{2}%
\renewcommand{\thetheorem}{3(a)}%
\begin{theorem}[闭环误差动态]\label{thm:3a}
误差坐标 $(e_z,e_\xi)$ 满足二阶系统
\begin{equation}
\dot e_z=A(\tilde x)\,e_\xi,\qquad
\dot e_\xi=-K_d e_\xi-A^{\top}(\tilde x)\,K_p\,e_z+d(t),
\tag{5.5}\label{eq:5.5}
\end{equation}
其中 $d=J\boldsymbol w_{\mathrm{dyn}}$，并按 \eqref{eq:5.1d} 分解为 $d=\Theta u_{\mathrm{fb}}+d_{\mathrm{ex}}$。
\end{theorem}

\begin{proof}
对 \eqref{eq:4.4} 逐项求导，用伴随求导法则 $\frac{d}{dt}\bigl(\mathrm{Ad}_{\tilde x}\boldsymbol a\bigr)=\mathrm{Ad}_{\tilde x}\dot{\boldsymbol a}+\mathrm{ad}_{\tilde{\boldsymbol\xi}}(\mathrm{Ad}_{\tilde x}\boldsymbol a)$（由 $\dot{\tilde x}=\tfrac12\tilde{\boldsymbol\xi}\tilde x$ 与其共轭直接展开）得 twist 误差动态
\begin{equation}
\dot{\tilde{\boldsymbol\xi}}=\dot{\boldsymbol\xi}-\mathrm{Ad}_{\tilde x}\dot{\boldsymbol\xi}_d-\mathrm{ad}_{\tilde{\boldsymbol\xi}}\bigl(\mathrm{Ad}_{\tilde x}\boldsymbol\xi_d\bigr)
\tag{5.5a}\label{eq:5.5a}
\end{equation}
将 \eqref{eq:5.2} 与 $\ddot{\boldsymbol q}=\ddot{\boldsymbol q}_{\mathrm{ref}}+\boldsymbol w_{\mathrm{dyn}}$ 代入 \eqref{eq:3.5}，再代入上式，前馈相消，余下 $d=J\boldsymbol w_{\mathrm{dyn}}$ 即得 \eqref{eq:5.5}。
\end{proof}

% ---- 定理 3(b) ----
\renewcommand{\thetheorem}{3(b)}%
\begin{theorem}[无扰：水平集不变性、渐近收敛与局部指数稳定]\label{thm:3b}
设 $d\equiv0$，$K_p=\mathrm{diag}(K_{p,O},K_{p,T})$ 对称正定（旋转/平移块）。取
\begin{equation}
0<c<c^{*}\triangleq\tfrac12\lambda_{\min}(K_{p,O}),
\qquad
\Omega_c\triangleq\{(e_z,e_\xi):V\le c\}.
\tag{5.5b}\label{eq:5.5b}
\end{equation}
则：\textbf{(i)} $\dot V=-e_\xi^\top K_de_\xi\le0$，故 $\Omega_c$ 紧且正向不变；\textbf{(ii)} 初值在 $\Omega_c$ 内且 $\tilde\eta(0)>0$ 时，全程有
\begin{equation}
\tilde\eta(t)\ \ge\ \eta_0\triangleq\sqrt{1-\frac{2c}{\lambda_{\min}(K_{p,O})}}>0 ,
\tag{5.5c}\label{eq:5.5c}
\end{equation}
从而 $A(\tilde x)$ 在 $\Omega_c$ 上一致可逆（$\det A=-\tfrac18\tilde\eta\le-\tfrac18\eta_0<0$）；\textbf{(iii)} $\Omega_c$ 内一切轨迹满足 $(e_z,e_\xi)\to(0,0)$；\textbf{(iv)} $(0,0)$ 是\textbf{局部指数稳定}的。
\end{theorem}

\begin{proof}
关键步骤是 $A^\top$-整形反馈使交叉项 $e_z^\top K_pAe_\xi$ 精确相消，得 $\dot V=-e_\xi^\top K_de_\xi$；LaSalle 定理与 $A$ 在 $\Omega_c$ 上的一致可逆性（附录 C.2）给出渐近收敛。详见附录 C.4。
\end{proof}

\begin{remark}
初始 $\tilde\eta<0$ 时可将 HDQ 误差整体翻号，不改变 $\tilde{\boldsymbol\xi}$；上述水平集保证后续无需切换。$\blacksquare$
\end{remark}

% ---- 定理 3(c) ----
\renewcommand{\thetheorem}{3(c)}%
\begin{theorem}[$L_2$ 扰动：H∞ 二次型/Schur 补判据与旋转/平移分量拆分]\label{thm:3c}
设 $d=d_{L_2}\in L_2$。

\textbf{证书参数}：$\kappa,\gamma_a>0$ 不进入控制律，认证 $d\to e_\xi$ 的增益 $\gamma^{*}=\gamma_a\sqrt\kappa$（量纲 s）；零初值下 $\|e_\xi\|_{L_2}\le\gamma^{*}\|d\|_{L_2}$。证书族内的阻尼反演见 3(c$'$)。

\textbf{(c-1) 合并判据}（$K_d$ 任意对称正定）：对给定存储函数与供给率，点态耗散不等式 $\dot V\le-\tfrac1{2\kappa}\|e_\xi\|^2+\tfrac{\gamma_a^2}2\|d\|^2$ 对任意 $(e_\xi,d)$ 成立的充要条件为
\begin{equation}
M\triangleq\begin{bmatrix}K_d-\tfrac1{2\kappa}I_6 & -\tfrac12 I_6\\[2pt] -\tfrac12 I_6 & \tfrac{\gamma_a^2}2 I_6\end{bmatrix}\succeq0
\;\overset{\text{Schur}}{\Longleftrightarrow}\;
K_d\succeq\tfrac12\bigl(\kappa^{-1}+\gamma_a^{-2}\bigr)I_6 ,
\tag{5.6a}\label{eq:5.6a}
\end{equation}
此时
\begin{equation}
\int_0^\infty\kappa^{-1}\|e_\xi\|^2dt\le\gamma_a^2\int_0^\infty\|d_{L_2}\|^2dt+2V(0),
\tag{5.6}\label{eq:5.6}
\end{equation}
即加速度层扰动到 twist 误差能量的 $L_2$ 增益 $\le\gamma_a\sqrt\kappa$（零初值时退化为纯增益界；此界只约束 $e_\xi$）。

\textbf{(c-2) 分量拆分判据}（$K_d=\mathrm{diag}(K_\omega,K_v)$ 块对角，且 $K_p=\mathrm{diag}(K_{p,O},\,k_{p,T}I_3)$，平移刚度块须各向同性，必要性见附录 C.3）：记 $e_\xi=[\tilde\omega;\tilde v]$、$d=[d_\omega;d_v]$，$V_\omega\triangleq\tfrac12\|\tilde\omega\|^2+\tfrac12\mathcal O^\top K_{p,O}\mathcal O$、$V_v\triangleq\tfrac12\|\tilde v\|^2+\tfrac{k_{p,T}}2\|\mathcal T\|^2$（$V=V_\omega+V_v$）。若
\begin{equation}
K_\omega\succeq\tfrac12\bigl(\kappa_\omega^{-1}+\gamma_\omega^{-2}\bigr)I_3,
\qquad
K_v\succeq\tfrac12\bigl(\kappa_v^{-1}+\gamma_v^{-2}\bigr)I_3,
\tag{5.6b}\label{eq:5.6b}
\end{equation}
则旋转/平移两分量\textbf{各自独立}满足
\begin{equation}
\int_0^\infty\kappa_\omega^{-1}\|\tilde\omega\|^2dt\le\gamma_\omega^2\int_0^\infty\|d_\omega\|^2dt+2V_\omega(0),
\quad
\int_0^\infty\kappa_v^{-1}\|\tilde v\|^2dt\le\gamma_v^2\int_0^\infty\|d_v\|^2dt+2V_v(0),
\tag{5.6$'$}\label{eq:5.6p}
\end{equation}
两组证书参数可独立指定。
\end{theorem}

\begin{proof}
含扰时交叉项仍精确相消，得 $\dot V=-e_\xi^\top K_de_\xi+e_\xi^\top d$；将性能目标写为 $[e_\xi;d]^\top M[e_\xi;d]\ge0$，取 Schur 补即得充要判据 \eqref{eq:5.6a}；分量解耦依赖两处代数恒零（$\mathcal T\times\mathcal T=0$、$\mathcal T\cdot(\mathcal T\times\tilde\omega)=0$）。详见附录 C.5。
\end{proof}

% ---- 定理 3(c')：最优证书与 γ 反演 ----
\renewcommand{\thetheorem}{3(c$'$)}%
\begin{theorem}[最优证书与目标增益反演：$\gamma^{*}$ 调参]\label{thm:3c2}
记认证 $L_2$ 增益 $\gamma^{*}\triangleq\gamma_a\sqrt{\kappa}$，$\lambda\triangleq\lambda_{\min}(K_d)>0$。则：
\textbf{(i)}（下界）证书族 \eqref{eq:5.6a} 的任一可行点都满足 $\gamma^{*}\ge1/\lambda$；
\textbf{(ii)}（可达）$(\kappa,\gamma_a)=(\lambda^{-1},\lambda^{-1/2})$ 可行且 $\gamma^{*}=1/\lambda$。
故证书族可认证的最优（最小）增益为
\begin{equation}
\gamma^{*}_{\mathrm{opt}}\;=\;\min_{(\kappa,\gamma_a)\ \mathrm{feasible}}\gamma_a\sqrt{\kappa}\;=\;\frac{1}{\lambda_{\min}(K_d)} ,
\tag{5.6c}\label{eq:5.6c}
\end{equation}
最优点为 $\kappa=\gamma_a^2=1/\lambda_{\min}(K_d)$。\textbf{在该证书族内}，认证目标 $\gamma^{*}$ 当且仅当 $K_d\succeq(1/\gamma^{*})I_6$，可行点为 $(\kappa,\gamma_a)=(\gamma^{*},\sqrt{\gamma^{*}})$。式 \eqref{eq:5.6b} 同理给出最优分量界 $\gamma^{*}_\omega=1/\lambda_{\min}(K_\omega)$、$\gamma^{*}_v=1/\lambda_{\min}(K_v)$。
\end{theorem}

\begin{proof}
(i)~由 AM--GM 不等式（各项为正），$\kappa^{-1}+\gamma_a^{-2}\ \ge\ 2\sqrt{\kappa^{-1}\gamma_a^{-2}}\ =\ \dfrac{2}{\gamma_a\sqrt{\kappa}}\ =\ \dfrac{2}{\gamma^{*}}$；代入判据 \eqref{eq:5.6a} 得 $\lambda\ \ge\ \tfrac12\bigl(\kappa^{-1}+\gamma_a^{-2}\bigr)\ \ge\ 1/\gamma^{*}$。
(ii)~在 $\kappa=\gamma_a^2=\lambda^{-1}$ 处 $\kappa^{-1}+\gamma_a^{-2}=2\lambda$，\eqref{eq:5.6a} 取等成立，证书增益 $\gamma_a\sqrt{\kappa}=1/\lambda$；且 $\kappa=\lambda^{-1}>1/(2\lambda)$，可行域非空。
反演规则的必要性即 (i)；充分性：若 $\lambda\ge1/\gamma^{*}$，则点 $(\kappa,\gamma_a)=(\gamma^{*},\sqrt{\gamma^{*}})$ 满足 $\kappa^{-1}+\gamma_a^{-2}=2/\gamma^{*}\le2\lambda$，可行，且其证书增益恰为 $\gamma_a\sqrt{\kappa}=\gamma^{*}$。
\end{proof}

\begin{remark}
当 $K_p=\mathrm{diag}(p_OI_3,p_TI_3)$、$K_d=\mathrm{diag}(K_\omega,K_v)$ 时，认证值等于线性化 $H_\infty$ 范数；一般 SPD 增益下仅为上界（附录 C.7）。$\blacksquare$
\end{remark}

% ---- 定理 3(d) ----
\renewcommand{\thetheorem}{3(d)}%
\begin{theorem}[$L_\infty$ 扰动：乘性分量分离与 twist 误差的均方极限界]\label{thm:3d}
设 $K_p=\mathrm{diag}(K_{p,O},K_{p,T})\succ0$，(A3) 成立，扰动按 \eqref{eq:5.1d} 分解，$D_{\mathrm{ex}}\triangleq\|d_{\mathrm{ex}}\|_{L_\infty}<\infty$。记\textbf{有效阻尼}
\begin{equation}
\lambda_{\mathrm{eff}}\triangleq\lambda_{\min}(K_d)-\alpha\,\lambda_{\max}(K_d)\;>\;0
\tag{5.7a}\label{eq:5.7a}
\end{equation}
设轨迹在 $[0,T_\Omega]$ 内留在 $\Omega_c$（\eqref{eq:5.5b}），记\textbf{等效扰动幅值}
\begin{equation}
D\triangleq D_{\mathrm{ex}}+\alpha\,\lambda_{\max}(K_p)\Bigl(1+\sqrt{\tfrac{2c}{\lambda_{\min}(K_{p,T})}}\Bigr)\sqrt{\tfrac{2c}{\lambda_{\min}(K_p)}} .
\tag{5.7b}\label{eq:5.7b}
\end{equation}
则 twist 误差满足有限时段均方界
\begin{equation}
\frac1T\int_0^T\|e_\xi(t)\|^2\,dt\;\le\;\frac{D^2}{\lambda_{\mathrm{eff}}^2}+\frac{2V(0)}{\lambda_{\mathrm{eff}}\,T},
\qquad 0<T\le T_\Omega,
\tag{5.7c}\label{eq:5.7c}
\end{equation}
若该水平集前提对所有 $t\ge0$ 成立，则进一步有
\begin{equation}
\limsup_{T\to\infty}\ \mathrm{RMS}_{[0,T]}(e_\xi)
\;\triangleq\;\limsup_{T\to\infty}\Bigl(\frac1T\int_0^T\|e_\xi\|^2dt\Bigr)^{1/2}
\;\le\;\frac{D}{\lambda_{\mathrm{eff}}}
\;\xrightarrow[\ \alpha\to0\ ]{}\;\frac{\|d_{\mathrm{ex}}\|_{L_\infty}}{\lambda_{\min}(K_d)} .
\tag{5.7}\label{eq:5.7}
\end{equation}
这是以留域为前提的 twist 均方界，不证明受扰水平集不变性、全状态 ISS 或位姿稳态标度律。
\end{theorem}

\begin{proof}
将乘性扰动 $\Theta u_{\mathrm{fb}}$ 拆为阻尼折减项与等效扰动项，分别用 Young 不等式回收，积分即得 RMS 界 \eqref{eq:5.7}。详见附录 C.6。
\end{proof}

% ---------------------------------------------------------------------
% 第 6 章 仿真验证
% ---------------------------------------------------------------------
\section{仿真验证}

\subsection{平台与任务}

在 CoppeliaSim 中以 7-DoF KUKA LBR4+ 力矩模式仿真（控制周期 5 ms），名义动力学取自 \cite{Gaz14}。$t=2.5\,\text{s}$ 时 0.25 kg 未建模负载刚性附着，构成持续偏差型扰动 $d_{\mathrm{ex}}$。任务为抓取--搬运--圆周跟踪（半径 0.06 m，标准角速度 1.0 rad/s，高速工况 2.5 rad/s）。对比三种控制器：\textbf{C1}（本文，式 \eqref{eq:5.2}）、\textbf{C2}（\cite{Ch20}，twist 误差与本文相同但位姿反馈取螺旋对数）、\textbf{C3}（\cite{P2} 经速度伺服桥接至加速度级）。三律共享同一轨迹、初始条件与名义模型，仅 $\ddot{\boldsymbol q}_{\mathrm{ref}}$ 计算分支不同；DC 刚度均配平至 80，闭环极点 $\{-4,-20\}$，以隔离结构差异。

\subsection{验证目标：从定理到可观测预言}

仿真不追求“性能竞赛”式的百分比优胜，而是检验 §3--5 理论框架的三类可证伪预言：

\textbf{(V1) 静态刚度标度律 \eqref{eq:5.9}}：若定理 3(d) 的扰动模型与 §5.4 的线性化正确，则同一物理扰动 $d_{\mathrm{ex}}$ 应由不同增益档的稳态残差反演出一致幅值；且 1/4 旋转刚度折减因子应被独立确认。

\textbf{(V2) 证书分化的极限暴露}：C1 与 C2 共享解析前馈与几何一致 twist 误差，预期在近恒等域数值等价（确认“同信息集”）；但 C1 具备定理 3 的严格耗散等式与 H∞ 证书，C2/C3 没有。这一分化应在离开理想条件的极限工况（高速、噪声）下表现为 C1/C2 与 C3 的结构性分离——而非随机波动。

\textbf{(V3) 水平集工程余度}：定理 3(b) 给出水平集 $\Omega_c$ 的正向不变性，预期在实际运行中应留有充分余度，以确认局部结论的工程可用性。

\subsection{结果}

\subsubsection{静态刚度标度律与扰动反演（验证 V1）}

带载稳态（$t\ge 12.5\,\text{s}$）下，C1-tuned 与 C1-fast 的位姿残差如下：

\begin{center}
\begin{tabular}{@{}lccc@{}}
\toprule
档位 & $K_p=\mathrm{diag}(p_O I_3, p_T I_3)$ & $\|\mathcal T\|_{\mathrm{rms}}$ (m) & $\|\mathcal O\|_{\mathrm{rms}}$ \\
\midrule
tuned & $(320, 80)$ & $4.86\times10^{-3}$ & $4.27\times10^{-3}$ \\
fast  & $(720, 180)$ & $2.20\times10^{-3}$ & $1.93\times10^{-3}$ \\
\bottomrule
\end{tabular}
\end{center}

由 \eqref{eq:5.9} 反演等效扰动：
\begin{equation*}
\|d_v\| = k_{p,T}\|\mathcal T\|_{\mathrm{ss}}, \qquad \|d_\omega\| = \tfrac12\lambda(K_{p,O})\|\mathcal O\|_{\mathrm{ss}}.
\end{equation*}
得 tuned 档 $(\|d_v\|, \|d_\omega\|) = (0.389, 0.683)$，fast 档 $(0.396, 0.696)$，两档一致至 \textbf{1.93\%}。该强一致性（两个独立增益档反演出同一物理量）直接验证了定理 3(d) 扰动模型的结构正确性——若 1/4 旋转折减因子缺失，旋转反演值将与平移分量相差 3.5 倍，一致性立即破坏。

作为对照，base 档（$K_p=16I_6$，未补偿 1/4 折减，有效旋转刚度仅 4）带载姿态误差放大至 tuned 档的 12.3 倍，反演扰动偏离 35\%--38\%。这从反面验证了 §5.4 的 1/4 折减补偿规则的必要性：不补偿则旋转/平移带宽严重失配，\eqref{eq:5.9} 失效。

\subsubsection{极限工况下的证书分化（验证 V2）}

标准工况下，三律位置级残差差异 $<0.1\%$（静态刚度锁定），速度级 $\|e_\xi\|_{\mathrm{rms}}$（$\times10^{-3}$）如下：

\begin{center}
\begin{tabular}{@{}lccc@{}}
\toprule
工况 & C1 & C2 & C3 \\
\midrule
标准 & 0.940 & 0.940 & 0.955 \\
高速 ($\omega=2.5$) & 2.314 & 2.314 & 2.361 \\
噪声 & 3.164 & 3.165 & 3.234 \\
\bottomrule
\end{tabular}
\end{center}

\textbf{近恒等域的结构等价}：C1 与 C2 在全部工况下数值等价（相对差 $\le 0.05\%$），确认二者共享同一信息集（解析前馈 + $\mathrm{Ad}$-搬运 twist 误差）。差异仅在于位姿反馈整形（$A^\top$ vs 螺旋对数），在近恒等域为高阶项——这定量支持了本文的核心主张：C1 相对 C2 的优势不在精度，而在证书（定理 3 的耗散等式、水平集不变性、均方界对 $A^\top$ 整形成立，对螺旋对数不成立）。

\textbf{极限工况的结构性分离}：C3 在噪声下劣化最明显（3.234 vs 3.165，相对 +2.2\%），与其桥接差分 $\Delta\dot q_{\mathrm{cmd}}/\Delta t$ 的噪声放大机制一致；在高速下劣化 +2.0\%，源于其一阶前馈无法解析补偿向心加速度的 $\omega^2$ 项。C1/C2 的 TNDQ 解析前馈避免了数值差分，故在极限工况下保持紧致——这验证了“有严格动力学证书”与“无证书桥接”在工程上的分化。

\subsubsection{水平集余度（验证 V3）}

杯附着后负载突变瞬间，$V^{\mathrm{base}}$ 峰值 $2.47\times10^{-2}$，在 $1.50\,\text{s}$ 内回落至 $3.36\times10^{-4}$；换算至 tuned 档权重后 $V^{\mathrm{tuned}}_{\mathrm{peak}}\le 0.494$，距定理 3(b) 水平集阈值 $c^*=160$ 余度约 \textbf{2.5 个数量级}。这表明定理 3(b) 的局部结论在实际运行中具有充足的工程裕度，并非仅存在于任意小邻域的理论陈述。

\subsection{讨论：理论优势的仿真映射}

上述结果可映射回 §3--5 的理论结构如下：

\textbf{(i) 定理 3(d) 的实用性}：$L_\infty$ 均方界 \eqref{eq:5.7} 预测了偏差型扰动下的稳态误差水平，而 \eqref{eq:5.9} 进一步给出了增益整定的解析规则——仿真中 1.93\% 的反演一致性确认了该规则的有效性。

\textbf{(ii) 定理 3(b) 的工程余度}：水平集 $\Omega_c$ 在实际运行中留有 2.5 个数量级余度，说明局部指数稳定结论在典型工况下是保守但可用的。

\textbf{(iii) 证书分化的可观测性}：C1 与 C2 的数值等价确认了“同信息集”前提，使得“C1 有证书、C2 无证书”成为可辩护的理论分化而非精度吹嘘；C3 在极限工况下的系统性劣化（噪声/高速）则从反面验证了动力学证书的价值——无严格前馈结构时，离散化与差分误差在应力下累积。

\textbf{(iv) 局限}：每组仅单次运行，1.5\%--2\% 的相对差异不具备统计显著性；真机实验与更大负载范围为后续工作。

% ---------------------------------------------------------------------
% 第 7 章 结论
% ---------------------------------------------------------------------
\section{结论}

本文基于三项代数 $\mathcal A_2$（TODQ）重构机械臂运动学：串联链正运动学一次 $O(n)$ 连乘同时输出位姿、twist 与加速度；误差体系定义在两项 HDQ 元素上，一次乘法同时生成右不变位姿误差与几何一致 twist 误差（定理 1），并闭合为严格级联误差运动学（定理 2）。几何一致计算力矩律使闭环满足精确耗散等式，在同一存储函数上给出三类证书：无扰时的水平集不变性与局部指数稳定（定理 3(b)）、$L_2$ 扰动下 H∞ 增益的 Schur 补判据（定理 3(c)）、$L_\infty$ 扰动下 twist 误差的均方极限界（定理 3(d)）。加速度层被证明无需误差项——期望加速度走前馈、不确定性走扰动通道，误差层仅需 DQ/HDQ 乘法。近恒等线性化 \eqref{eq:5.8} 还揭示了一个实现陷阱：旋转刚度存在 1/4 折减，不补偿则旋转/平移带宽失配。

CoppeliaSim 力矩模式仿真（§6）定量核验了上述主张：标度律 \eqref{eq:5.9} 的扰动反演在两个独立增益档上一致至 1.93\%，base 档 12.3 倍的姿态误差放大从反面确认折减补偿的必要性；本文控制律与二阶基线 \cite{Ch20} 数值等价（$\le 0.05\%$），在高速、噪声工况下优于一阶桥接基线（$+2.0\%$--$2.2\%$）；水平集余度约 2.5 个数量级，证书条件全程未被违反。与 \cite{Ch20} 的数值等价表明，本文的贡献是同性能下的严格证书与大误差几何鲁棒性，而非精度提升。

\textbf{局限与后续工作}：strictification（附录 C.4）未完成，$\Omega_c$ 前提尚未自洽闭合；稳定性结论为工作域局部（unwinding 与 $\tilde\eta=0$ 处 $A$ 的奇异为拓扑障碍）；真机实验待验证。


% ================================ 附录 ================================
% ---------------------------------------------------------------------
% 附录 A：DQ/HDQ 层的验证性推导
% 附录 C：控制层的补充推导
% （原文档即无附录 B 与 C.5，此处保持一致）
% ---------------------------------------------------------------------

\section*{附录 A：DQ/HDQ 层的验证性推导}

\subsection*{A.1 定理 1 的 (i)(ii)(iii)}

(ii)：由 \eqref{eq:4.3} 直接右乘 $\tilde x$。(i)：翻转不变性：$2(-\dot{\tilde x})(-\tilde x^*)=2\dot{\tilde x}\tilde x^*$。(iii)：$\dot{\tilde x}=\dot{\hat{\underline x}}\hat{\underline x}_d^{\,*}+\hat{\underline x}\dot{\hat{\underline x}}_d^{\,*}$；用 $\dot{\hat{\underline x}}=\tfrac12\boldsymbol\xi\hat{\underline x}$ 与 $\dot{\hat{\underline x}}_d^{\,*}=-\tfrac12\hat{\underline x}_d^{\,*}\boldsymbol\xi_d$（后者由 $\dot{\hat{\underline x}}_d=\tfrac12\boldsymbol\xi_d\hat{\underline x}_d$ 取共轭）：
$\dot{\tilde x}=\tfrac12\boldsymbol\xi\tilde x-\tfrac12\hat{\underline x}\hat{\underline x}_d^{\,*}\boldsymbol\xi_d=\tfrac12\boldsymbol\xi\tilde x-\tfrac12\tilde x\boldsymbol\xi_d$。右乘 $2\tilde x^*$：$\tilde{\boldsymbol\xi}=\boldsymbol\xi-\tilde x\boldsymbol\xi_d\tilde x^*=\boldsymbol\xi-\mathrm{Ad}_{\tilde x}\boldsymbol\xi_d$。含速度级扰动 $\boldsymbol v_w,\boldsymbol v_c$（\cite{P2} 模型）时右端加 $\boldsymbol v_w+\boldsymbol v_c$。

\subsection*{A.2 引理 1 的证明}

\begin{equation*}
\tfrac{d}{dt}(\tilde x\boldsymbol a\tilde x^*)
=\dot{\tilde x}\boldsymbol a\tilde x^*+\tilde x\dot{\boldsymbol a}\tilde x^*+\tilde x\boldsymbol a\dot{\tilde x}^* .
\end{equation*}
代入 $\dot{\tilde x}=\tfrac12\tilde{\boldsymbol\xi}\tilde x$、$\dot{\tilde x}^*=-\tfrac12\tilde x^*\tilde{\boldsymbol\xi}$：
\begin{equation*}
=\tfrac12\tilde{\boldsymbol\xi}(\mathrm{Ad}_{\tilde x}\boldsymbol a)+\mathrm{Ad}_{\tilde x}\dot{\boldsymbol a}-\tfrac12(\mathrm{Ad}_{\tilde x}\boldsymbol a)\tilde{\boldsymbol\xi}
=\mathrm{Ad}_{\tilde x}\dot{\boldsymbol a}+\mathrm{ad}_{\tilde{\boldsymbol\xi}}(\mathrm{Ad}_{\tilde x}\boldsymbol a).
\end{equation*}
\eqref{eq:5.4}：对 \eqref{eq:4.4} 逐项求导并取 $\boldsymbol a=\boldsymbol\xi_d$。$\blacksquare$

\section*{附录 C：控制层的补充推导}

\subsection*{C.1 扰动项的适定性（\eqref{eq:5.1} 为何是显式等式，以及 $\alpha$ 条件的真正作用）}

$\boldsymbol\tau=\hat M\ddot{\boldsymbol q}_{\mathrm{ref}}+\hat C\dot{\boldsymbol q}+\hat g$ 代入真实动力学 $M\ddot{\boldsymbol q}+C\dot{\boldsymbol q}+\boldsymbol g=\boldsymbol\tau+\delta\boldsymbol\tau+\boldsymbol\tau_{\mathrm{ext}}$ 并解出 $\ddot{\boldsymbol q}$：
\begin{equation*}
\ddot{\boldsymbol q}=M^{-1}\hat M\ddot{\boldsymbol q}_{\mathrm{ref}}+M^{-1}\bigl(\Delta C\dot{\boldsymbol q}+\Delta\boldsymbol g+\delta\boldsymbol\tau+\boldsymbol\tau_{\mathrm{ext}}\bigr),
\qquad M^{-1}\hat M=I+M^{-1}\Delta M ,
\end{equation*}
即 \eqref{eq:5.1}。注意 $\ddot{\boldsymbol q}_{\mathrm{ref}}$ 由 \eqref{eq:5.2} 完全由\textbf{当前可测量} $(\boldsymbol q,\dot{\boldsymbol q},\hat{\underline x}_d,\boldsymbol\xi_d,\dot{\boldsymbol\xi}_d)$ 给出，\textbf{不依赖 $\ddot{\boldsymbol q}$}，故 \eqref{eq:5.1} 对 $\ddot{\boldsymbol q}$ 是\textbf{显式赋值}：不存在隐式代数环，既不需要移项、也不需要 Neumann 级数，更不会产生 $\alpha/(1-\alpha)$ 型的等效扰动增益因子。

由此可见 $\alpha<1/\mathrm{cond}_2(K_d)$（\eqref{eq:5.1f}）与适定性无关，而是为了使定理 3(d) 中的\textbf{有效阻尼保持正定}：乘性分量 $\Theta u_{\mathrm{fb}}$ 中与 $-K_de_\xi$ 同向的部分会削弱阻尼，由 $|e_\xi^\top\Theta(-K_de_\xi)|\le\alpha\lambda_{\max}(K_d)\|e_\xi\|^2$ 得 $\lambda_{\mathrm{eff}}=\lambda_{\min}(K_d)-\alpha\lambda_{\max}(K_d)$，正性恰好等价于 \eqref{eq:5.1f}。$\alpha$ 条件是一个小增益型的证书条件，而非微分方程适定性条件；若 $\alpha$ 过大，闭环仍然定义良好，但本文的 Lyapunov 证书失效。

\subsection*{C.2 定理 3(b) 的定量奇异值下界与收敛论证细节}

$\dot V=-e_\xi^\top K_de_\xi$ 只是关于全状态的\textbf{半}负定（$e_\xi=0$ 面上 $\dot V=0$），故不能直接给出收敛，必须补 LaSalle 步骤（定理 3(b) 证明第 4 步），而 LaSalle 步骤需要 $A$ 在工作域内可逆。定量化如下。

由 \eqref{eq:4.5}，$A=\begin{bmatrix}A_{11}&0\\ -[\mathcal T]_\times& I_3\end{bmatrix}$，$A_{11}=-\tfrac12(\tilde\eta I+[\mathcal O]_\times)$。其逆为
\begin{equation*}
\begin{gathered}
A^{-1}=\begin{bmatrix}A_{11}^{-1}&0\\ [\mathcal T]_\times A_{11}^{-1}& I_3\end{bmatrix},\\[2pt]
A_{11}^{-1}=-2\,\frac{\tilde\eta^2I+\tilde\eta[\mathcal O]_\times^{\top}+\mathcal O\mathcal O^\top}{\tilde\eta(\tilde\eta^2+\|\mathcal O\|^2)}
=-\frac2{\tilde\eta}\bigl(\tilde\eta^2I-\tilde\eta[\mathcal O]_\times+\mathcal O\mathcal O^\top\bigr),
\end{gathered}
\end{equation*}
用到 $\tilde\eta^2+\|\mathcal O\|^2=1$。$\|A_{11}^{-1}\|_2$ 无需放缩即可精确算出：记 $B\triangleq\tilde\eta I+[\mathcal O]_\times$（故 $A_{11}=-\tfrac12B$），则
\begin{equation*}
B^\top B=(\tilde\eta I-[\mathcal O]_\times)(\tilde\eta I+[\mathcal O]_\times)=\tilde\eta^2I-[\mathcal O]_\times^2=(\tilde\eta^2+\|\mathcal O\|^2)I-\mathcal O\mathcal O^\top=I-\mathcal O\mathcal O^\top ,
\end{equation*}
其特征值为 $\{\,\tilde\eta^2,\,1,\,1\,\}$（沿 $\mathcal O$ 方向为 $1-\|\mathcal O\|^2=\tilde\eta^2$，两个正交方向为 1）。于是 $B$ 的奇异值为 $\{\tilde\eta,1,1\}$，同时给出两个\textbf{精确}等式
\begin{equation*}
\|A_{11}\|_2=\tfrac12\sigma_{\max}(B)=\tfrac12 ,
\quad
\|A_{11}^{-1}\|_2=\frac2{\sigma_{\min}(B)}=\frac2{\tilde\eta} ,
\end{equation*}
前者正是 §5.2 与定理 3(b) 反复使用的精确值。再由分块下三角结构 $\|A^{-1}\|_2\le(1+\|\mathcal T\|)\allowbreak\|A_{11}^{-1}\|_2+1$，
\begin{equation*}
\sigma_{\min}(A)=\frac1{\|A^{-1}\|_2}\;\ge\;\Bigl[\frac{2(1+\|\mathcal T\|)}{\tilde\eta}+1\Bigr]^{-1} .
\end{equation*}
在 $\Omega_c$ 内 $\tilde\eta\ge\eta_0$（\eqref{eq:5.5c}）且 $\|\mathcal T\|\le\sqrt{2c/\lambda_{\min}(K_{p,T})}$，故 $\sigma_{\min}(A)\ge c_A(\eta_0,c)>0$ 一致成立。于是 LaSalle 第 4 步中由 $A^\top K_pe_z\equiv0$ 推出 $\|e_z\|\le\|K_p^{-1}\|_2\|A^{-\top}\|_2\cdot0=0$ 严格成立。局部指数率由线化矩阵 $F$ 的谱给出，逐分量具体极点见 \eqref{eq:5.8}。域限制 $\tilde\eta>0$ 与 \cite{P2} Remark 1 的 unwinding 条件一致。

\begin{remark}
注：无扰情形下 LaSalle 路线自身是完整的，但不能改用级联 ISS 论证闭合。
\end{remark}

\subsection*{C.3 定理 3(c) 的二次型/Schur 补细节}

\textbf{(c-1) Schur 补判据。} 记供给率 $s(e_\xi,d)\triangleq-\tfrac1{2\kappa}\|e_\xi\|^2+\tfrac{\gamma_a^2}2\|d\|^2$。性能目标 $\dot V\le s$ 等价于 $-e_\xi^\top K_de_\xi+e_\xi^\top d-s(e_\xi,d)\le0$，整理为 $-[e_\xi;d]^\top M[e_\xi;d]\le0$，其中 $M$ 为 \eqref{eq:5.6a} 的分块矩阵。$M\succeq0$ 且右下块 $\tfrac{\gamma_a^2}2I\succ0$，取 Schur 补得
\begin{equation*}
M\succeq0\iff K_d-\tfrac1{2\kappa}I-\tfrac1{2\gamma_a^2}I\succeq0\iff K_d\succeq\tfrac12(\kappa^{-1}+\gamma_a^{-2})I .
\end{equation*}
这是\textbf{当且仅当}判据：不定号交叉项 $e_\xi^\top d$ 保留在二次型内整体判定，无任何符号放缩。

\textbf{(c-2) 两处恒零的几何含义与失效条件。} 拆分精确成立依赖两条叉积混合积恒等式：(i) $(A^\top e_z)_\omega$ 中 $[\mathcal T]_\times\mathcal T=\mathcal T\times\mathcal T=0$；(ii) $\dot{\mathcal T}=-[\mathcal T]_\times\tilde\omega+\tilde v$ 中耦合项对 $\tfrac12\|\mathcal T\|^2$ 不做功，$\mathcal T^\top(\mathcal T\times\tilde\omega)=0$。二者是代数恒等式，不依赖工作域，故 (c-2) 与 (c-1) 一样是全局结论。若 $K_d$ 非块对角，$-e_\xi^\top K_de_\xi$ 含 $\tilde\omega^\top K_{\omega v}\tilde v$ 型交叉项，两分量能量不再分离，退回合并判据 (c-1)。

\textbf{$K_{p,T}$ 各向同性的必要性。} 将平移刚度写作一般对称正定 $K_{p,T}$ 时，上述两处恒零均失效：(i) $(A^\top K_pe_z)_\omega=A_{11}^\top K_{p,O}\mathcal O+[\mathcal T]_\times K_{p,T}\mathcal T$，而 $[\mathcal T]_\times K_{p,T}\mathcal T=\mathcal T\times(K_{p,T}\mathcal T)\ne0$ 除非 $\mathcal T$ 是 $K_{p,T}$ 的特征向量；(ii) 平移储能变为 $\tfrac12\mathcal T^\top K_{p,T}\mathcal T$，其导数中的耦合项 $-\mathcal T^\top K_{p,T}[\mathcal T]_\times\tilde\omega$ 不再为零（仅当 $K_{p,T}=k_{p,T}I_3$ 时由 $\mathcal T^\top[\mathcal T]_\times=0$ 而消失）。两项残留均是 $\tilde\omega$--$\mathcal T$ 型耦合，使 $V_\omega,V_v$ 不再各自闭合，退回 (c-1)。

\subsection*{C.4 strictification 与局部 ISS（未来工作）}

定理 3(d) 的三条缺口的共同根源是 $\dot V$ 中没有 $-\|e_z\|^2$ 型负项。标准补救是 strictification：取 $W=V+\epsilon e_z^\top K_pA(\tilde x)e_\xi$（$\epsilon>0$ 待定），在 $\Omega_c$ 内 $\epsilon$ 足够小时 $W$ 与 $V$ 等价，而求导后新增 $-\epsilon e_z^\top K_pAA^\top K_pe_z$ 项提供 $e_z$ 方向负定。若能验证剩余交叉项（含 $\dot A$ 项）可被两个负定项吸收，则 $\dot W\le-c_1\|(e_z,e_\xi)\|^2+c_2\|d\|\cdot\|(e_z,e_\xi)\|$，得到真正的局部 ISS-Lyapunov 函数。本文未完成此验证（关键难点是 $\dot A$ 项的一致界与 $\epsilon$ 的可行区间非空性），列为 §7 局限 (i)。

\subsection*{C.6 定理 3(b) 的完整证明}

\begin{enumerate}[leftmargin=2em]
  \item \textbf{交叉项精确相消}：沿 \eqref{eq:5.5}（$d\equiv0$），
  \begin{equation*}
  \dot V=e_\xi^\top\dot e_\xi+e_z^\top K_p\dot e_z
  =e_\xi^\top\bigl(-K_de_\xi-A^\top K_pe_z\bigr)+e_z^\top K_pAe_\xi
  =-e_\xi^\top K_de_\xi ,
  \end{equation*}
  因 $e_z^\top K_pAe_\xi=(A^\top K_p^\top e_z)^\top e_\xi=(A^\top K_pe_z)^\top e_\xi$——这里\textbf{只}用到 $K_p$ 对称与 $K_p$ 写在 $A^\top$ 内侧，不需要 $K_p$ 为标量。故 $\dot V\le0$，$\Omega_c$ 正向不变；又 $V\le c$ 给出 $\|e_\xi\|\le\sqrt{2c}$、$\|e_z\|\le\sqrt{2c/\lambda_{\min}(K_p)}$，故 $\Omega_c$ 紧。

  \item \textbf{工作域保号 \eqref{eq:5.5c}}：$V\le c$ 蕴含 $\tfrac12\mathcal O^\top K_{p,O}\mathcal O\le c$，即 $\|\mathcal O\|^2\le2c/\lambda_{\min}(K_{p,O})<1$（由 $c<c^*$）。由 $\mathcal O=-\mathrm{Im}\,\tilde r$ 与 $\|\tilde r\|=1$ 得 $\tilde\eta^2+\|\mathcal O\|^2=1$，故 $|\tilde\eta|\ge\eta_0>0$；$\tilde\eta(t)$ 连续且恒不为零，符号不可突变，由 $\tilde\eta(0)>0$ 得 \eqref{eq:5.5c}。

  \item \textbf{$A$ 的行列式}：$A$ 块下三角（\eqref{eq:4.5}），故
  \begin{equation*}
  \det A=\det\bigl(-\tfrac12(\tilde\eta I_3+[\mathcal O]_\times)\bigr)\cdot\det I_3
  =\bigl(-\tfrac12\bigr)^3\tilde\eta\bigl(\tilde\eta^2+\|\mathcal O\|^2\bigr)=-\tfrac18\tilde\eta ,
  \end{equation*}
  用到 $\det(aI_3+[b]_\times)=a(a^2+\|b\|^2)$ 与 $\tilde\eta^2+\|\mathcal O\|^2=1$。定量奇异值下界 $\sigma_{\min}(A)\ge\bigl[2(1+\|\mathcal T\|)/\tilde\eta+1\bigr]^{-1}$ 见附录 C.2。

  \item \textbf{LaSalle}：$\Omega_c$ 紧且不变，$E\triangleq\{\dot V=0\}\cap\Omega_c=\{e_\xi=0\}\cap\Omega_c$。若轨迹全程留在 $E$ 内：$e_\xi\equiv0\Rightarrow\dot e_\xi\equiv0\Rightarrow A^\top K_pe_z\equiv0$；由第 3 步 $A$ 可逆、$K_p\succ0$ 得 $e_z\equiv0$。故 $E$ 内最大不变集为 $\{(0,0)\}$，由 LaSalle 不变集定理（\cite{Kha02} Thm 4.4）得 (iii)。

  \item \textbf{局部指数稳定}：在 $(0,0)$ 处 $\tilde x\to1$，$A\to A_0=\mathrm{diag}(-\tfrac12I_3,I_3)$，\eqref{eq:5.5} 的雅可比为
  \begin{equation*}
  F=\begin{bmatrix}0_6 & A_0\\ -A_0^\top K_p & -K_d\end{bmatrix}.
  \end{equation*}
  对该 LTI 系统同一个 $V$ 仍给出 $\dot V=-e_\xi^\top K_de_\xi\le0$，且 $A_0$ 可逆，重复第 4 步得 LTI 系统渐近稳定，故 $F$ 为 Hurwitz；再由 Lyapunov 线化定理（\cite{Kha02} Thm 4.7）得非线性系统在原点邻域指数稳定。块对角 $K_d,K_p$ 下 $F$ 逐分量解耦，其两个二阶多项式与极点由 \eqref{eq:5.8} 显式给出。
\end{enumerate}

\subsection*{C.7 定理 3(c) 的完整证明}

\textbf{第一步（$\dot V$ 精确式）}：含扰时定理 3(b) 证明第 1 步的交叉项相消与 $d$ 无关，故
\begin{equation*}
\dot V=-e_\xi^\top K_de_\xi+e_\xi^\top d
\end{equation*}
精确成立（此式不含任何放缩；$e_\xi^\top d$ 可正可负）。

\textbf{第二步（(c-1) 判据的等价性）}：性能目标即 $-e_\xi^\top K_de_\xi+e_\xi^\top d+\tfrac1{2\kappa}\|e_\xi\|^2-\tfrac{\gamma_a^2}2\|d\|^2\le0$，等价于二次型不等式
\begin{equation*}
\begin{bmatrix}e_\xi\\ d\end{bmatrix}^{\!\top}\!M\!
\begin{bmatrix}e_\xi\\ d\end{bmatrix}\ge0\quad\forall(e_\xi,d)\in\mathbb R^{12},
\end{equation*}
即 $M\succeq0$。右下块 $\tfrac{\gamma_a^2}2I\succ0$，取 Schur 补得 $K_d-\tfrac1{2\kappa}I-\tfrac1{2\gamma_a^2}I\succeq0$，即 \eqref{eq:5.6a}。关键在于不定号交叉项 $e_\xi^\top d$ 保留在二次型内整体判定，不经任何符号放缩。

\textbf{第三步（全局存在性）}：由 $M\succeq0$ 得 $\dot V\le\tfrac{\gamma_a^2}2\|d\|^2$，故 $V(t)\le V(0)+\tfrac{\gamma_a^2}2\|d_{L_2}\|_{L_2}^2<\infty$，$(e_z,e_\xi)$ 一致有界，解在 $[0,\infty)$ 上存在（无有限时间逃逸）。

\textbf{第四步（积分收尾）}：在 $[0,T]$ 上积分 $\dot V\le-\tfrac1{2\kappa}\|e_\xi\|^2+\tfrac{\gamma_a^2}2\|d\|^2$，弃去 $V(T)\ge0$，令 $T\to\infty$（单调收敛）即得 \eqref{eq:5.6}。

\textbf{第五步（(c-2) 分量解耦）}：块对角 $K_d$ 与各向同性平移刚度 $K_{p,T}=k_{p,T}I_3$ 下两分量储能精确解耦，关键是两处混合积恒零：由 \eqref{eq:4.5}，$(A^\top K_pe_z)_\omega=A_{11}^\top K_{p,O}\mathcal O+k_{p,T}[\mathcal T]_\times\mathcal T=A_{11}^\top K_{p,O}\mathcal O$（$\mathcal T\times\mathcal T=0$，故旋转反馈不含 $\mathcal T$）、$(A^\top K_pe_z)_v=k_{p,T}\mathcal T$；又 $\dot{\mathcal T}=-[\mathcal T]_\times\tilde\omega+\tilde v$ 中的耦合项做功为零（$\mathcal T\cdot(\mathcal T\times\tilde\omega)=0$）。于是位姿交叉项在两分量内分别由 $K_{p,O}$ 对称与 $k_{p,T}$ 为标量而精确相消，
\begin{equation*}
\dot V_\omega=-\tilde\omega^\top K_\omega\tilde\omega+\tilde\omega^\top d_\omega,
\qquad
\dot V_v=-\tilde v^\top K_v\tilde v+\tilde v^\top d_v ,
\end{equation*}
对每个分量重复第二至第四步的论证（$I_6\to I_3$）即得 \eqref{eq:5.6b}$\Rightarrow$\eqref{eq:5.6p}。两处恒零为代数恒等式，故 (c-1)/(c-2) 的全部结论均不依赖工作域 $\tilde\eta>0$。

\subsection*{C.8 定理 3(d) 的完整证明}

\begin{enumerate}[leftmargin=2em]
  \item \textbf{精确耗散等式与乘法项展开}：由定理 3(c) 证明第一步，$\dot V=-e_\xi^\top K_de_\xi+e_\xi^\top d$ 精确成立（无任何放缩）。代入 \eqref{eq:5.1d} 与 \eqref{eq:5.2} 的 $u_{\mathrm{fb}}=-K_de_\xi-A^\top K_pe_z$：
  \begin{equation*}
  \dot V=-e_\xi^\top K_de_\xi\;\underbrace{-\,e_\xi^\top\Theta K_de_\xi}_{\text{与阻尼同类}}\;\underbrace{-\,e_\xi^\top\Theta A^\top K_pe_z}_{\text{与扰动同类}}\;+\;e_\xi^\top d_{\mathrm{ex}} .
  \end{equation*}

  \item \textbf{两类乘法项的分别回收}：$|e_\xi^\top\Theta K_de_\xi|\le\alpha\lambda_{\max}(K_d)\|e_\xi\|^2$（回收进阻尼，得 \eqref{eq:5.7a} 的 $\lambda_{\mathrm{eff}}$；由 (A3)/\eqref{eq:5.1f} 即 $\alpha<\lambda_{\min}(K_d)/\lambda_{\max}(K_d)$ 得 $\lambda_{\mathrm{eff}}>0$）；$|e_\xi^\top\Theta A^\top K_pe_z|\allowbreak\le\alpha\|A\|_2\allowbreak\lambda_{\max}(K_p)\|e_z\|\,\|e_\xi\|$（回收进等效扰动幅值 $D$）。

  \item \textbf{$A$ 的谱范数界}：由 $A_{11}^\top A_{11}=\tfrac14(I_3-\mathcal O\mathcal O^\top)$ 得 $\|A_{11}\|_2=\tfrac12$（\textbf{精确值}，与 $\tilde x$ 无关；奇异值计算见附录 C.2），再由 $A$ 的块下三角结构得 $\|A\|_2\le\max\{\|A_{11}\|_2,1\}+\|[\mathcal T]_\times\|_2=1+\|\mathcal T\|$。在 $\Omega_c$ 上 $\|\mathcal T\|\le\sqrt{2c/\lambda_{\min}(K_{p,T})}$、$\|e_z\|\le\sqrt{2c/\lambda_{\min}(K_p)}$，代入第 2 步即得 \eqref{eq:5.7b} 与
  \begin{equation}
  \dot V\ \le\ -\lambda_{\mathrm{eff}}\|e_\xi\|^2+D\,\|e_\xi\| .
  \tag{5.7d}\label{eq:5.7d}
  \end{equation}

  \item \textbf{Young 与积分收尾}：$D\|e_\xi\|\le\tfrac{\lambda_{\mathrm{eff}}}2\|e_\xi\|^2+\tfrac{D^2}{2\lambda_{\mathrm{eff}}}$，故 $\dot V\le-\tfrac{\lambda_{\mathrm{eff}}}2\|e_\xi\|^2+\tfrac{D^2}{2\lambda_{\mathrm{eff}}}$。在 $[0,T]$ 上积分并弃去 $V(T)\ge0$：
  \begin{equation*}
  \begin{gathered}
  \tfrac{\lambda_{\mathrm{eff}}}2\int_0^T\|e_\xi\|^2dt\\[2pt]
  \ \le\ V(0)+\tfrac{D^2}{2\lambda_{\mathrm{eff}}}\,T ,
  \end{gathered}
  \end{equation*}
  两端除以 $\tfrac{\lambda_{\mathrm{eff}}}2T$ 即 \eqref{eq:5.7c}；令 $T\to\infty$ 得 \eqref{eq:5.7}。$\alpha\to0$ 时 $D\to D_{\mathrm{ex}}$、$\lambda_{\mathrm{eff}}\to\lambda_{\min}(K_d)$，退化为经典形式。
\end{enumerate}


% ================================ 参考文献 ================================
% 编号与标签均与原文档一致（[P1]、[P2]、[Ch20] 等；第 3--5 条原文无标签，
% 按其在列表中的序号 [3]、[4]、[5] 引用）。
\begin{thebibliography}{18}
\raggedright

\bibitem[P1]{P1}
A.~Cohen, M.~Shoham, \emph{Hyper Dual Quaternions representation of rigid bodies kinematics},
Mechanism and Machine Theory 150 (2020) 103861.

\bibitem[P2]{P2}
L.F.C. Figueredo, B.V. Adorno, J.Y. Ishihara, \emph{Robust H$\infty$ kinematic control of manipulator robots using dual quaternion algebra},
Automatica 132 (2021) 109817.

\bibitem[3]{Fike11}
J.A. Fike, J.J. Alonso, \emph{The Development of Hyper-Dual Numbers for Exact Second-Derivative Calculations}, AIAA 2011-886.

\bibitem[4]{Lynch17}
K.M. Lynch, F.C. Park, \emph{Modern Robotics: Mechanics, Planning, and Control}, Cambridge University Press, 2017.

\bibitem[5]{Khatib87}
O. Khatib, \emph{A unified approach for motion and force control of robot manipulators: The operational space formulation},
IEEE J. Robotics and Automation 3(1), 1987.

\bibitem[Spo92]{Spo92}
M.W. Spong, \emph{On the robust control of robot manipulators},
IEEE Trans. Automatic Control 37(11), 1992.

\bibitem[Kha02]{Kha02}
H.K. Khalil, \emph{Nonlinear Systems}, 3rd ed., Prentice Hall, 2002.

\bibitem[Ch20]{Ch20}
A. Chandra, J.A. Corrales-Ramon, Y. Mezouar, \emph{Resolved-acceleration control of serial robotic manipulators using unit dual quaternions},
IFAC-PapersOnLine 53(2) (2020) 8500--8505.

\bibitem[Gaz14]{Gaz14}
C. Gaz, F. Flacco, A. De Luca, \emph{Identifying the dynamic model used by the KUKA LWR: A reverse engineering approach},
Proc. IEEE Int. Conf. Robotics and Automation (ICRA), 2014, 1386--1392.

\bibitem[Roh13]{Roh13}
E. Rohmer, S.P.N. Singh, M. Freese, \emph{V-REP: A versatile and scalable robot simulation framework},
Proc. IEEE/RSJ Int. Conf. Intelligent Robots and Systems (IROS), 2013, 1321--1326.

\bibitem[LWP80]{LWP80}
J.Y.S. Luh, M.W. Walker, R.P.C. Paul, \emph{Resolved-acceleration control of mechanical manipulators},
IEEE Trans. Automatic Control 25(3), 1980.

\bibitem[Abd91]{Abd91}
C. Abdallah, D. Dawson, P. Dorato, M. Jamshidi, \emph{Survey of robust control for rigid robots},
IEEE Control Systems Magazine 11(2), 1991.

\bibitem[Sag99]{Sag99}
H.G. Sage, M.F. De Mathelin, E. Ostertag, \emph{Robust control of robot manipulators: a survey},
International Journal of Control 72(16), 1999.

\bibitem[ZDG96]{ZDG96}
K. Zhou, J.C. Doyle, K. Glover, \emph{Robust and Optimal Control}, Prentice Hall, 1996.

\bibitem[Nak86]{Nak86}
Y. Nakamura, H. Hanafusa, \emph{Inverse kinematic solutions with singularity robustness for robot manipulator control},
ASME J. Dynamic Systems, Measurement, and Control 108(3), 1986.

\bibitem[Nak08]{Nak08}
J. Nakanishi, R. Cory, M. Mistry, J. Peters, S. Schaal, \emph{Operational space control: A theoretical and empirical comparison},
International Journal of Robotics Research 27(6), 2008.

\bibitem[Ber93]{Ber93}
H. Berghuis, H. Nijmeijer, \emph{A passivity approach to controller--observer design for robots},
IEEE Trans. Robotics and Automation 9(6), 1993.

\bibitem[Con25]{Con25}
D. Condurache, \emph{An overview of higher-order kinematics of rigid body and multibody systems with nilpotent algebra},
Mechanism and Machine Theory 209 (2025) 105959.

\end{thebibliography}

\end{document}
